
\documentclass{article}
\usepackage{colortbl}
\usepackage{makecell}
\usepackage{multirow}
\usepackage{supertabular}

\begin{document}

\newcounter{utterance}

\centering \large Interaction Transcript for game `wordle\_withclue', experiment `medium\_frequency\_words\_clue\_no\_critic', episode 14 with deepseek{-}v3.2{-}t1.0.
\vspace{24pt}

{ \footnotesize  \setcounter{utterance}{1}
\setlength{\tabcolsep}{0pt}
\begin{supertabular}{c@{$\;$}|p{.15\linewidth}@{}p{.15\linewidth}p{.15\linewidth}p{.15\linewidth}p{.15\linewidth}p{.15\linewidth}}
    \# & \multicolumn{2}{c}{Player} && \multicolumn{2}{c}{Game Master} \\
    \hline

    \theutterance \stepcounter{utterance}  
    & & & \multicolumn{4}{p{0.6\linewidth}}{
        \cellcolor[rgb]{0.9,0.9,0.9}{
            \makecell[{{p{\linewidth}}}]{
                \texttt{\tiny{[P1$\langle$GM]}}
                \texttt{You are a language wizard who likes to guess words by using the given rules.} \\
\\ 
\texttt{Welcome to Wordle! You have six attempts to guess the target word, a valid English word of five lowercase letters (a{-}z). Please use the tags "explanation:" and "guess:" to provide a concise explanation for each guess.} \\
\\ 
\texttt{To help you make an informed guess, you will receive a clue for the word, such as} \\
\texttt{clue: snowy white.} \\
\\ 
\texttt{Here is an example guess based on the clue:} \\
\texttt{explanation: In the fairy tail Snow White, the girl is killed because she eats a poisoned apple. And the word apple has 5 letters.} \\
\texttt{guess: apple} \\
\\ 
\texttt{After each guess, your answer will be validated, and you will receive feedback indicating which letters are correct (green), which letters are correct but in the wrong position (yellow), and which letters are incorrect (red). This feedback can be useful in determining which letters to include or exclude in your next guess.} \\
\\ 
\texttt{For example, the feedback for "apple" might be:} \\
\texttt{guess\_feedback: a<yellow> p<yellow> p<green> l<yellow> e<red>} \\
\\ 
\texttt{The explanation should contain details about how both the guess\_feedback and clue work together in tandem to formulate a new guess.} \\
\\ 
\texttt{Do not generate any other text.} \\
\\ 
\texttt{Let's begin with your first guess.} \\
\\ 
\texttt{clue: nine{-}to{-}five routine} \\
            }
        }
    }
     \\ \\

    \theutterance \stepcounter{utterance}  
    & \multicolumn{4}{p{0.6\linewidth}}{
        \cellcolor[rgb]{0.561,0.737,0.561}{
            \makecell[{{p{\linewidth}}}]{
                \texttt{\tiny{[P1$\rangle$GM]}}
                \texttt{explanation: The clue "nine{-}to{-}five routine" strongly suggests a word related to daily work or schedule, likely "daily". This is a common five{-}letter word fitting the clue.} \\
\texttt{guess: daily} \\
            }
        }
    }
    & & \\ \\

    \theutterance \stepcounter{utterance}  
    & & & \multicolumn{2}{p{0.3\linewidth}}{
        \cellcolor[rgb]{0.9,0.9,0.9}{
            \makecell[{{p{\linewidth}}}]{
                \texttt{\tiny{[GM$|$GM]}}
                \texttt{attempts = 1} \\
\texttt{target\_word = grind} \\
\texttt{target\_word\_clue = nine{-}to{-}five routine} \\
\texttt{guess = daily} \\
\texttt{guess\_feedback = d<yellow> a<red> i<green> l<red> y<red>} \\
            }
        }
    }
    & & \\ \\

    \theutterance \stepcounter{utterance}  
    & & & \multicolumn{4}{p{0.6\linewidth}}{
        \cellcolor[rgb]{0.9,0.9,0.9}{
            \makecell[{{p{\linewidth}}}]{
                \texttt{\tiny{[P1$\langle$GM]}}
                \texttt{guess\_feedback: d<yellow> a<red> i<green> l<red> y<red>} \\
\\ 
\texttt{Provide your response only in this format.} \\
\texttt{explanation: details} \\
\texttt{guess: word} \\
            }
        }
    }
     \\ \\

    \theutterance \stepcounter{utterance}  
    & \multicolumn{4}{p{0.6\linewidth}}{
        \cellcolor[rgb]{0.561,0.737,0.561}{
            \makecell[{{p{\linewidth}}}]{
                \texttt{\tiny{[P1$\rangle$GM]}}
                \texttt{explanation: Feedback shows 'i' is correct and green in position 3. 'd' is present but wrong spot (yellow). 'a', 'l', 'y' are absent (red). The clue "nine{-}to{-}five routine" still suggests a work{-}related word. With 'i' fixed in the 3rd spot and 'd' present somewhere else, possible words include "drill" (as in routine practice) or "dried". "Drill" fits both the routine concept and the letter pattern (d r i ? ?). Since 'l' is red, the last two letters must not be 'll'. Trying "drift" but that doesn't fit the clue. "Droid" doesn't fit. Reconsidering: "idler" is not a routine. "Dirge" is not routine. "Drink" is not routine. Perhaps "diner" (as in a lunch routine)? But 'a' and 'l' are red, so 'a' is not in it. "Diver" doesn't fit. Wait, 'i' is green at position 3, pattern is \_ \_ i \_ \_. 'd' is yellow, so not in position 1. Possible pattern: ? ? i ? ? with 'd' in positions 2,4, or 5. Considering common five{-}letter words for routine: "habit" fits routine but has 'a' which is red. "Chore" has no 'i'. "Shift" fits nine{-}to{-}five (work shift) perfectly. Check letters: s h i f t. 'i' is green at position 3. 'd' is yellow, but 'shift' has no 'd'. However, feedback said 'd' is yellow, so it must be included. But if 'd' is required, maybe the word is "idiot"? No. Maybe the routine is "drive" (as in commute)? But 'a' and 'l' are red, 'drive' has no 'a' or 'l', but it has 'v' and 'e'. Feedback: d<yellow> a<red> i<green> l<red> y<red>. So letters a,l,y are out. 'd' is in but not first. 'i' third. So possible words: "dizzy", but 'y' is red, so no. "ditch", 't' and 'h' possible. "ditto" fits pattern? d i t t o, but 'i' is in position 2, not 3. So not. Let's think of common routines: "dairy" but 'a' and 'y' are red. "diary" but 'a' and 'y' are red. New approach: The clue "nine{-}to{-}five routine" could be "shift". But 'd' is yellow, so shift doesn't have 'd'. Maybe the word is "daily" but feedback ruled out a,l,y. Perhaps "duty" is four letters. Could it be "ditto"? Not routine. Maybe "drill" but 'l' is red. What about "droll"? 'l' red, 'o' and double 'l'? 'l' red so no 'l'. Given the feedback, the word likely starts with a letter other than d, contains 'd' not in first, has 'i' third, and excludes a,l,y. Let's guess a word that fits the clue and this pattern: "grind" (as in daily grind). Check: g r i n d. 'i' is third (green). 'd' is yellow (present but not in fifth position? Actually feedback says d<yellow> from first guess where d was first. In "grind", d is fifth, which could be yellow if not in correct spot. But we need to see if d's correct spot is known. It was yellow from first guess, meaning in "daily", d was first but correct spot is not first. So in target, d is not first. In "grind", d is fifth, which could be correct or not. Also, letters g,r,n are new and not yet tested. This seems plausible. "Grind" means routine work. So I'll guess "grind".} \\
\texttt{guess: grind} \\
            }
        }
    }
    & & \\ \\

    \theutterance \stepcounter{utterance}  
    & & & \multicolumn{2}{p{0.3\linewidth}}{
        \cellcolor[rgb]{0.9,0.9,0.9}{
            \makecell[{{p{\linewidth}}}]{
                \texttt{\tiny{[GM$|$GM]}}
                \texttt{attempts = 2} \\
\texttt{target\_word = grind} \\
\texttt{target\_word\_clue = nine{-}to{-}five routine} \\
\texttt{guess = grind} \\
\texttt{guess\_feedback = g<green> r<green> i<green> n<green> d<green>} \\
            }
        }
    }
    & & \\ \\

    \theutterance \stepcounter{utterance}  
    & & & \multicolumn{2}{p{0.3\linewidth}}{
        \cellcolor[rgb]{0.9,0.9,0.9}{
            \makecell[{{p{\linewidth}}}]{
                \texttt{\tiny{[GM$|$GM]}}
                \texttt{game\_result = WIN} \\
            }
        }
    }
    & & \\ \\

\end{supertabular}
}

\end{document}
